\exercisesfor{6}

The conventions of § 6 remain valid unless otherwise mentioned.

¶ 1. Let g be a semi-simple Lie algebra and ρ a finite-dimensional representation of g on M.

\begin{enumerate}
\def\labelenumi{(\alph{enumi})}
\item
  If \(\rho\) is simple and non-zero, \(\mathrm{H}^p (\mathfrak{g},\mathbf{M}) = \{0\}\) for all \(p\) . (Use Proposition 1 and Exercise 12 (j) of §3.)
\item
  For all \(\rho\) , \(\mathbf{H}^1(\mathfrak{g}, \mathbf{M}) = \{0\}\) . (If \(\rho\) is simple and non-zero, apply (a). If \(\rho\) is zero, use the fact that \(\mathfrak{g} = \mathcal{D}\mathfrak{g}\) . In the general case, argue by induction on the dimension of \(\rho\) : if N is a sub-g-module of M distinct from \(\{0\}\) and M, use the exact sequence:
\end{enumerate}

\[
\mathrm{H} ^ {1} (\mathfrak {g}, \mathrm{N}) \rightarrow \mathrm{H} ^ {1} (\mathfrak {g}, \mathrm{M}) \rightarrow \mathrm{H} ^ {1} (\mathfrak {g}, \mathrm{M} / \mathrm{N})
\]

established in Exercise 12 (f) of § 3.) Hence recover Remark 2 of no. 2.

\begin{enumerate}
\def\labelenumi{(\alph{enumi})}
\setcounter{enumi}{2}
\item
  Deduce from (b) a proof of Theorem 2. (Use Exercise 12 (g) of §3.) Deduce also from (b) that every derivation of g is inner. (Use Exercise 12 (h) of §3.)
\item
  For all \(\rho\) , \(\mathrm{H}^2(\mathfrak{g}, \mathrm{M}) = \{0\}\) . (Arguing as for (b), it suffices to consider the case where \(\rho = 0\) . Let \(c \in \mathrm{H}^2(\mathfrak{g}, \mathrm{M})\) . Consider, following Exercise 12 (i) of § 3, the central extension \(\mathfrak{h}\) of \(\mathfrak{g}\) by \(\mathrm{M}\) defined by \(c\) . The adjoint representation of \(\mathfrak{h}\) defines a representation of \(\mathfrak{g}\) on \(\mathfrak{h}\) . By Theorem 2, \(\mathrm{M}\) admits in \(\mathfrak{h}\) a supplement which is stable under \(\mathfrak{g}\) and hence the extension is trivial. Hence \(c = 0\) by Exercise 12 (i) of § 3.)
\item
  Deduce from (b) and (d) a proof of Theorem 5. (As in the text, it can be reduced to the case where the radical is commutative.)
\end{enumerate}

\begin{enumerate}
\def\labelenumi{\arabic{enumi}.}
\setcounter{enumi}{1}
\item
  Let g be a Lie algebra, r its radical and \((a_{0}, a_{1}, \ldots)\) a sequence of ideals of g defined as follows: (1) \(a_{0} = \{0\}\) ; (2) \(a_{i+1}/a_{i}\) is a maximal commutative ideal of \(g/a_{i}\) . Let p be the smallest integer such that \(a_{p} = a_{p+1} = \cdots\) . Show that \(r = a_{p}\) . (Show that \(g/a_{p}\) is semi-simple.)
\item
  For a Lie algebra g to be semi-simple, it is necessary and sufficient that it be reductive in every Lie algebra containing g as a subalgebra. (If g satisfies this condition, let M be a g-module of finite dimension over K. Considering M as a commutative algebra, form the semi-direct product of g and M, in which g is reductive. Deduce that M is semi-simple.)
\item
  Let g be a semi-simple Lie algebra and \(\rho\) a non-zero simple representation of g on a finite-dimensional space M. Let h be the corresponding semi-direct product. Show that \(h = Dh\) , that the centre of h is zero and that h is not the product of a semi-simple algebra and a solvable algebra.
\item
  Let \(\mathfrak{a}\) be a Lie algebra and \(\mathfrak{r}\) its radical. If \(\mathfrak{r}\) has a decreasing sequence of characteristic ideals \(\mathfrak{r} = \mathfrak{r}_0 \supset \mathfrak{r}_1 \supset \dots \supset \mathfrak{r}_n = \{0\}\) such that \(\dim \mathfrak{r}_i / \mathfrak{r}_{i+1} = 1\) for \(0 \leqslant i < n\) , \(\mathfrak{g}\) is the product of a semi-simple algebra and a solvable algebra.
\end{enumerate}

(Let s be a Levi subalgebra of g. For all \(x \in s\) , let \(\rho(x)\) be the restriction of \(ad_{g}x\) to r. Then \(\rho\) is a direct sum of 1-dimensional representations, which are zero because \(s = Ds\) .)

\begin{enumerate}
\def\labelenumi{\arabic{enumi}.}
\setcounter{enumi}{5}
\item
  Let \(\mathfrak{g}\) be a Lie algebra and \(\mathcal{D}^{\infty}\mathfrak{g}\) the intersection of the \(\mathcal{D}^{p}\mathfrak{g}\) ( \(p = 1,2,\ldots\) ). Show that \(\mathfrak{g}/\mathcal{D}^{\infty}\mathfrak{g}\) is solvable and that \(\mathcal{D}(\mathcal{D}^{\infty}\mathfrak{g}) = \mathcal{D}^{\infty}\mathfrak{g}\) . For \(\mathfrak{g}\) to be isomorphic to the product of a semi-simple algebra and a solvable algebra, it is necessary and sufficient that \(\mathcal{D}^{\infty}\mathfrak{g}\) be semi-simple.
\item
  \begin{enumerate}
  \def\labelenumii{(\alph{enumii})}
  \tightlist
  \item
    Let g be a Lie algebra, b a semi-simple ideal of g and a the centralizer of b in g, so that g is identified with \(b \times a\) . Show that, for every ideal \(\ell\) of g, then \(\ell = (\ell \cap \mathfrak{b}) \times (\ell \cap \mathfrak{a})\) . (Let \(\ell_{1}\) be the canonical projection of \(\ell\) onto b; it is an ideal of b and hence \(\mathcal{D}\ell_{1} = \ell_{1}\) ; deduce that \(\ell_{1} \subset \ell\) .)
  \end{enumerate}
\end{enumerate}

\begin{enumerate}
\def\labelenumi{(\alph{enumi})}
\setcounter{enumi}{1}
\item
  Let \(\beta\) be an invariant bilinear form on \(\mathfrak{g}\) . Show that \(\mathfrak{b}\) and \(\mathfrak{a}\) are orthogonal with respect to \(\beta\) (use the fact that \(\mathfrak{b} = \mathcal{D}\mathfrak{b}\) and that \(\beta = \beta_{1} + \beta_{2}\) , where \(\beta_{1}\) (resp. \(\beta_{2}\) ) is an invariant bilinear form whose restriction to \(\mathfrak{a}\) (resp. \(\mathfrak{b}\) ) is zero.
\item
  Deduce from (a) that there exists in \(\mathfrak{g}\) a largest semi-simple ideal. (Consider a maximal semi-simple ideal of \(\mathfrak{g}\) .)
\end{enumerate}

\begin{enumerate}
\def\labelenumi{\arabic{enumi}.}
\setcounter{enumi}{7}
\tightlist
\item
  Let g be a Lie algebra. An ideal h of g is called minimal if \(h \neq \{0\}\) and every ideal of g contained in h is equal to \(\{0\}\) or h.
\end{enumerate}

\begin{enumerate}
\def\labelenumi{(\alph{enumi})}
\item
  Every simple ideal of g is minimal.
\item
  Let \(\mathfrak{h}\) be a minimal ideal of \(\mathfrak{g}\) and \(\mathfrak{r}\) the radical of \(\mathfrak{g}\) . Then either \(\mathfrak{h} \subset \mathfrak{r}\) , in which case \(\mathfrak{h}\) is abelian, or \(\mathfrak{h} \cap \mathfrak{r} = \{0\}\) , in which case \(\mathfrak{h}\) is simple. (Use the fact that the derived ideal of a Lie algebra and the simple components of a semi-simple Lie algebra are characteristic ideals.)
\end{enumerate}

\begin{enumerate}
\def\labelenumi{\arabic{enumi}.}
\setcounter{enumi}{8}
\item
  For a Lie algebra g to be reductive, it is necessary and sufficient that its centre c be equal to its largest nilpotent ideal. (If the condition holds, let r be the radical of g; Dr is contained in the centre of r, hence r is nilpotent and hence r = c.)
\item
  Let \(\mathfrak{g}\) be a Lie algebra such that the conditions \(x \in \mathfrak{g}, y \in \mathfrak{g}, [[x, y], y] = 0\) imply \([x, y] = 0\) . Show that \(\mathfrak{g}\) is reductive. (Show that the radical \(\mathfrak{r}\) of \(\mathfrak{g}\) is commutative using Exercise 3 of §5. Then show that \([\mathfrak{g}, \mathfrak{r}] = 0\) .)
\end{enumerate}

¶ 11. (a) Let g be a Lie algebra, r its radical, s a Levi subalgebra of g and m an ideal of g containing r. There exists an ideal t of s, supplementary to m in g, such that \([m, t] \subset r\).

\begin{enumerate}
\def\labelenumi{(\alph{enumi})}
\setcounter{enumi}{1}
\tightlist
\item
  Let g be a Lie algebra and a subinvariant subalgebra of g. Then there exists a composition series \(g = g_{0} \supset g_{1} \supset \cdots \supset g_{k} = a\) such that \(g_{i}\) is the direct sum of \(g_{i+1}\) and a subalgebra \(b_{i}\) which is either 1-dimensional and contained in the radical \(r(g_{i})\) of \(g_{i}\) or simple and such that \([b_{i}, g_{i+1}] \subset r(g_{i+1})\) . (Reduce it to the case where a is an ideal of g. Let \(g' = g/a\) and \(r'\) be the radical of \(g'\) . The algebra \(g'/r'\) is the product of its simple ideals \(a_{1}, a_{2}, \ldots, a_{c}\) .
\end{enumerate}

Take a composition series of \(\mathfrak{g}'\) consisting of a composition series of \(\mathbf{r}'\) whose successive quotients are 1-dimensional and the inverse images of the ideals \(a_1, a_1 \times a_2, \ldots, a_1 \times a_2 \times \cdots \times a_c\) . Then take the inverse image in \(\mathfrak{g}\) of this composition series.)

\begin{enumerate}
\def\labelenumi{\arabic{enumi}.}
\setcounter{enumi}{11}
\tightlist
\item
  Let g be a Lie algebra, r its radical and D a derivation of g.
\end{enumerate}

\begin{enumerate}
\def\labelenumi{(\alph{enumi})}
\tightlist
\item
  If \(\mathrm{D}(\mathfrak{r}) = \{0\} \), D is inner. (Let c be the centralizer of r in g. The adjoint representation of g defines a representation \(x^{*} \mapsto \rho(x^{*})\) of \(g^{*} = g/r\) on c.~On the other hand, \([\mathrm{D}(\mathfrak{g}), \mathfrak{r}] \subset \mathrm{D}([\mathfrak{g}, \mathfrak{r}]) + [\mathfrak{g}, \mathrm{D}(\mathfrak{r})] = \{0\}\) , hence \(\mathrm{D}(\mathfrak{g}) \subset c\) and hence D defines a linear mapping \(D^{*}: g^{*} \to c\) . Show that
\end{enumerate}

\[
\mathrm{D} ^ {*} ([ x ^ {*}, y ^ {*} ]) = \rho (x ^ {*}) \mathrm{D} ^ {*} y ^ {*} - \rho (y ^ {*}) \mathrm{D} ^ {*} x ^ {*}
\]

for all \(x^{*}, y^{*}\) in \(\mathfrak{g}^{*}\) . Deduce that there exists \(a \in \mathfrak{c}\) such that \(D^{*}x^{*} = \rho(x^{*})a\) for all \(x^{*} \in \mathfrak{g}^{*}\) (cf.~no. 2, Remark 2), whence \(Dx = [x, a]\) for all \(x \in \mathfrak{g}\) .)

\begin{enumerate}
\def\labelenumi{(\alph{enumi})}
\setcounter{enumi}{1}
\tightlist
\item
  If D coincides on r with an inner derivation of g, D is inner. (Apply (a).)
\end{enumerate}

\begin{enumerate}
\def\labelenumi{\arabic{enumi}.}
\setcounter{enumi}{12}
\item
  Let g be a Lie algebra over an algebraically closed field, r its radical and \(\rho\) a finite-dimensional simple representation of g. Then \(\rho(z)\) is scalar for \(z \in r\) . (Reduce it to the case where g is reductive. In this case, r is the centre of g.)
\item
  Let X be a nilpotent matrix in \(\mathfrak{gl}(n, \mathbf{K})\) . Show that there exist in \(\mathfrak{gl}(n, \mathbf{K})\) two other matrices Y, H such that \([H, X] = X\) , \([H, Y] = -Y\) , \([X, Y] = H\) . (Argue by induction on n, using the Jordan form of X; thus reduce it to the case where \(X = E_{12} + E_{23} + \cdots + E_{n-1,n}\) ; then take Y to be a combination of the \(E_{k,k-1}\) and H a combination of the \(E_{jj'}\) .)
\end{enumerate}

¶ 15. (a) Let X, Y be two matrices of \(\mathfrak{gl}(n, \mathrm{K})\) . Show that, if \([X, Y] = X\) , X is nilpotent. (For every polynomial \(f \in K[T]\) , observe that

\[
[ f (X), Y ] = X f ^ {\prime} (X).)
\]

\begin{enumerate}
\def\labelenumi{(\alph{enumi})}
\setcounter{enumi}{1}
\item
  Let \(\mathfrak{g}\) be a Lie algebra and \(h\) , \(x\) two elements of \(\mathfrak{g}\) such that \([h, x] = x\) and there exists \(z \in \mathfrak{g}\) with \([x, z] = h\) . Show that there exists \(y \in \mathfrak{g}\) such that \([x, y] = h\) and \([h, y] = -y\) . (Observe first that \([z, h] + z\) belongs to the centralizer \(\mathfrak{n}\) of \(x\) in \(\mathfrak{g}\) . Then show that \(\mathfrak{n}\) is stable under ad \(h\) and that the restriction to \(\mathfrak{n}\) of I - ad \(h\) is bijective; to do this, note that ad \(x\) is nilpotent, using \((a)\) , and prove that, if \(\mathfrak{g}_{\mathfrak{l}}\) is the image of \(\mathfrak{g}\) under \((\mathrm{ad}x)^{\mathfrak{l}}\) , I - ad \(h\) gives on taking quotients a bijection of \((\mathfrak{n} \cap \mathfrak{g}_{\mathfrak{l}-1}) / (\mathfrak{n} \cap \mathfrak{g}_{\mathfrak{l}})\) ; for this use the relation \([\mathrm{ad}z, (\mathrm{ad}x)^{\mathfrak{k}}] = -k(\mathrm{ad}x)^{\mathfrak{k}-1}(\mathrm{ad}h + \frac{k-1}{2}I).)\)
\item
  Let \(\mathfrak{g}\) be a subalgebra of \(\mathfrak{gl}(n, \mathbf{K})\) ; suppose that, for every nilpotent matrix \(X \in \mathfrak{g}\) , there exist two other matrices \(H, Y\) in \(\mathfrak{g}\) such that \([H, X] = X\) , \([H, Y] = -Y\) , \([X, Y] = H\) . Let \(\mathfrak{h}\) be a subalgebra of \(\mathfrak{g}\) such that there exists
\end{enumerate}

a subspace m supplementary to h in g and such that \([h, m] \subset m\) . Show that h has the same property as g for all its nilpotent matrices (use (b)).

\begin{enumerate}
\def\labelenumi{\arabic{enumi}.}
\setcounter{enumi}{15}
\item
  Let g be a subalgebra of \(\mathfrak{gl}(n,\mathbf{K})\) such that the identity representation of g is semi-simple. Show that every nilpotent matrix \(X\in g\) is contained in a 3-dimensional simple Lie subalgebra of g. (Show that g is reductive in \(\mathfrak{gl}(n,\mathbf{K})\) ; then use Exercises 14 and 15 (c).)
\item
  Show that the algebra derived from a real simple Lie algebra by extending the base field from \(\mathbf{R}\) to \(\mathbf{C}\) is not always simple. (Let \(\mathfrak{g}\) be a real simple Lie algebra. If \(\mathfrak{g}' = \mathfrak{g}_{(\mathbf{C})}\) is simple, let \(\mathfrak{h}\) be the real Lie algebra derived from \(\mathfrak{g}'\) by restricting the field of scalars to \(\mathbf{R}\) . We know that \(\mathfrak{h}\) is simple. Then \(\mathfrak{h}_{(\mathbf{C})}\) is, by Exercise 4 of §1, the product of two algebras isomorphic to \(\mathfrak{g}'\) . Cf. also Exercise 26 (b).
\item
  \begin{enumerate}
  \def\labelenumii{(\alph{enumii})}
  \tightlist
  \item
    Let g be a simple Lie algebra. Every invariant bilinear form on g is either zero or non-degenerate. If K is algebraically closed, every invariant bilinear form \(\beta\) on g is proportional to the Killing form \(\beta_{0}\) . (Consider the endomorphism \(\sigma\) of the vector space g defined by \(\beta(x,y)=\beta_{0}(\sigma(x),y)\) and show that \(\sigma\) commutes with \(ad_{g}x\) and is hence scalar.) Show that this result may be false if K is not algebraically closed. (Use Exercise 17 and the fact that the dimension of the space of invariant bilinear forms does not change when the base field is extended.)
  \end{enumerate}
\end{enumerate}

\begin{enumerate}
\def\labelenumi{(\alph{enumi})}
\setcounter{enumi}{1}
\item
  Let g be a semi-simple Lie algebra over an algebraically closed field. Deduce from (a) and Exercise 7 (b) that the dimension of the space of invariant bilinear forms on g is equal to the number of simple components of g and that all these forms are symmetric.
\item
  Let g be a simple Lie algebra, M the dual space of the vector space g, with the dual representation of the adjoint representation, and b the semidirect product of g and M defined by \(x \mapsto x_{M}\) . For y, z in g and \(y'\) , \(z'\) in M, we write \(\beta(y + y', z + z') = \langle y, z' \rangle + \langle z, y' \rangle\) . Show that \(\beta\) is a non-degenerate invariant symmetric bilinear form on b, which is associated with no representation of b. (Observe that the radical and nilpotent radical of b are equal to M.)
\end{enumerate}

\begin{enumerate}
\def\labelenumi{\arabic{enumi}.}
\setcounter{enumi}{18}
\tightlist
\item
  Let \(g\) be a reductive Lie algebra, \(U\) its enveloping algebra, \(Z\) the centre of \(U\) and \(V\) the subspace of \(U\) generated by the elements of the form \(uv - vu\) ( \(u \in U, v \in U\) ).
\end{enumerate}

\begin{enumerate}
\def\labelenumi{(\alph{enumi})}
\item
  U is the direct sum of Z and V. (Apply Proposition 6 of §3 to the representation \(x \mapsto \mathrm{ad}_U x\) of \(\mathfrak{g}\) on U, observing that, in the decomposition \(\mathrm{U} = \sum_{n} \mathrm{U}^n\) of §2, no. 7, Corollary 4 to Theorem 1, the \(\mathrm{U}^n\) are stable under this representation.)
\item
  Let \(u \mapsto u^{\natural}\) be the projector of U onto Z parallel to V. Show that \((uv)^{\natural} = (vu)^{\natural}, (zu)^{\natural} = zu^{\natural}\) for all \(u \in \mathrm{U}, v \in \mathrm{U}, z \in \mathrm{Z}\) .
\item
  Let R be a two-sided ideal of U. Show that \(\mathbf{R} = (\mathbf{R} \cap \mathbf{Z}) + (\mathbf{R} \cap \mathbf{V})\) . (To show that the component in V of an element r of R belongs to R, reduce it to the case where K is algebraically closed; observe that R is stable under the representation \(\rho\) of U on U which extends \(x \mapsto ad_{U}x\) ; decompose r into its components in the \(U^{n}\) ; finally, apply Corollary 2 to Theorem 1 of Algebra, Chapter VIII, § 4, no. 2 to the restriction of \(\rho\) to \(U^{n}\) .)
\item
  Let \(\mathbf{R}'\) be an ideal of \(Z\) . Let \(\mathbf{R}_1\) be the two-sided ideal of \(U\) generated by \(\mathbf{R}'\) . Show that \(\mathbf{R}_1 \cap Z = \mathbf{R}'\) . (Use (b).)
\end{enumerate}

\begin{enumerate}
\def\labelenumi{\arabic{enumi}.}
\setcounter{enumi}{19}
\tightlist
\item
  Suppose that \(\mathbf{K}\) is algebraically closed. Let \(\mathfrak{g}\) be a Lie algebra and \(\mathfrak{c}\) its centre.
\end{enumerate}

\begin{enumerate}
\def\labelenumi{(\alph{enumi})}
\item
  If \(\mathfrak{g}\) admits a finite-dimensional faithful simple representation, \(\mathfrak{g}\) is reductive and \(\dim \mathfrak{c} \leqslant 1\) .
\item
  Conversely, if g is reductive and dim \(c \leqslant 1\) , g admits a finite-dimensional faithful simple representation. (If g is simple, use the adjoint representation. If \(g = g_{1} \times g_{2}\) and \(g_{1}, g_{2}\) admit finite-dimensional faithful simple representations \(\rho_{1}, \rho_{2}\) on spaces \(M_{1}, M_{2}\) , show that
\end{enumerate}

\[
(x _ {1}, x _ {2}) \mapsto \rho_ {1} (x _ {1}) \otimes 1 + 1 \otimes \rho_ {2} (x _ {2})
\]

is a semi-simple representation \(\rho\) of \(g\) , then that the centralizer of the \(g\) -module \(\mathrm{M}_1 \otimes \mathrm{M}_2\) reduces to the scalars and hence that \(\rho\) is simple. When \(g_1\) and \(g_2\) are semi-simple or \(g_1\) is semi-simple and \(g_2\) is commutative, show that \(\rho\) is faithful by considering the kernel of \(\rho\) which is an ideal of \(g_1 \times g_2\) .

\begin{enumerate}
\def\labelenumi{\arabic{enumi}.}
\setcounter{enumi}{20}
\item
  Let V be a vector space over K of finite dimension n \textgreater{} 1. Show that \(\mathfrak{sl}(V)\) is simple. (Reduce it to the case where K is algebraically closed. Suppose that \(\mathfrak{sl}(V) = a \times b\) , with dim a = a \textgreater{} 0, dim b = b \textgreater{} 0. Let A (resp. B) be the associative algebra generated by 1 and a (resp. b). Then V can be considered as a simple \((\mathrm{A} \otimes_{\mathrm{K}} \mathrm{B})\) -module. Hence there exists an A-module P of finite dimension p over K and a B-module Q of finite dimension q over K such that V is \((\mathrm{A} \otimes_{\mathrm{K}} \mathrm{B})\) -isomorphic to \(P \otimes_{K} Q\) (Algebra, Chapter VIII, §7, no. 7, Proposition 8 and no. 4, Theorem 2); P and Q are faithful and hence \(a \leqslant p^{2}\) , \(b \leqslant q^{2}\) . On the other hand, \(a + b = n^{2} - 1\) and pq = n, whence \((p^{2} - 1)(q^{2} - 1) \leqslant 2\) , which is contradictory.)
\item
  \begin{enumerate}
  \def\labelenumii{(\alph{enumii})}
  \tightlist
  \item
    Let \(\mathfrak{g}\) be a nilpotent Lie algebra over an algebraically closed field \(\mathbf{K}\) of arbitrary characteristic. Let \(\mathfrak{z}\) be its centre, \(\rho\) a finite-dimensional representation of \(\mathfrak{g}\) and \(\beta\) the associated bilinear form. Show that \(\mathfrak{z}\cap \mathcal{D}\mathfrak{g}\) is orthogonal to \(\mathfrak{g}\) with respect to \(\beta\) . (Reduce it, using a Jordan-Hölder series, to the case where \(\rho\) is simple. Then note that, for \(z\in \mathfrak{g}\cap \mathcal{D}\mathfrak{g},\rho (z)\) is a scalar matrix with trace zero. Then use Exercise 22 (b) of §4.) Deduce that, if \(\beta\) is nondegenerate, \(\mathfrak{g}\) is commutative. (Use Exercise 3 (a) of §4.)
  \end{enumerate}
\end{enumerate}

\begin{enumerate}
\def\labelenumi{(\alph{enumi})}
\setcounter{enumi}{1}
\item
  Suppose that \(\mathbf{K}\) is of characteristic 2. Show that for the solvable Lie algebra \(\mathfrak{gl}(2,\mathbf{K}) = \mathfrak{g}\) the bilinear form associated with the identity representation is non-degenerate, but that the centre \(\mathfrak{z}\) is contained in \(\mathcal{D}\mathfrak{g}\) and is \(\neq \{0\}\) .
\item
  Let g be the 6-dimensional Lie algebra over K with basis \((a, b, c, d, e, f)\) and multiplication table \([a, b] = -[b, a] = d\) , \([a, c] = -[c, a] = e\) , \([b, c] = -[c, b] = f\) and the other brackets are 0. Let β be the bilinear form on g such that β \((c, d) = \beta(d, c) = 1\) , β \((a, f) = \beta(f, a) = 1\) , β \((b, e) = \beta(e, b) = -1\) and the other values of β at ordered pairs of the basis of g in question are 0. Show that β is invariant, g nilpotent, \(\mathfrak{z} = \mathcal{D}g \neq \{0\}\) and β non-degenerate.
\end{enumerate}

\begin{enumerate}
\def\labelenumi{\arabic{enumi}.}
\setcounter{enumi}{22}
\tightlist
\item
  Let \(\mathfrak{g}\) be a 3-dimensional non-commutative Lie algebra over a field \(\mathbf{K}\) of arbitrary characteristic.
\end{enumerate}

\begin{enumerate}
\def\labelenumi{(\alph{enumi})}
\item
  If \(\mathfrak{g}\) has centre \(\mathfrak{z} \neq \{0\}\) , then dim \(\mathfrak{z} = 1\) (§ 1, Exercise 2) and dim \(\mathcal{D}_{\mathfrak{g}} = 1\) . If \(\mathfrak{z} \neq \mathcal{D}_{\mathfrak{g}}, \mathfrak{g}\) is the product of \(\mathfrak{z}\) and the 2-dimensional non-commutative algebra. If \(\mathfrak{z} = \mathcal{D}_{\mathfrak{g}}, \mathfrak{g}\) is the 3-dimensional non-commutative nilpotent algebra (§ 4, Exercise 9 (a)).
\item
  If \(\mathfrak{z} = \{0\}\) but there exists in \(\mathfrak{g}\) a 2-dimensional commutative subalgebra, this subalgebra is unique and is equal to \(\mathcal{D}_{\mathfrak{g}}\) ; for all \(x \notin \mathcal{D}_{\mathfrak{g}}\) , the restriction \(u\) of \(ad x\) to \(\mathcal{D}_{\mathfrak{g}}\) is a bijection of this vector space, determined to within a multiplicative constant. Conversely, every automorphism \(u\) of a 2-dimensional vector space \(\mathrm{Ka} + \mathrm{Kb}\) determines a Lie algebra structure on \(\mathfrak{g} = \mathrm{Ka} + \mathrm{Kb} + \mathrm{Kc}\) by the conditions \([a, b] = 0\) , \([c, a] = u(a)\) , \([c, b] = u(b)\) . For two Lie algebras thus defined by automorphisms \(u_1, u_2\) to be isomorphic, it is necessary and sufficient that the matrices of \(u_1\) and \(u_2\) be similar to within a scalar factor.
\item
  Suppose that there exists in g no commutative subalgebra of dimension \textgreater1. Show that there then exists \(a \in g\) such that \(a \notin (\operatorname{ad} a)g\) (assuming that an element x does not have this property, show, using the Jacobi identity, that another element has the desired property). There then exists a basis \((a, b, c)\) of g such that \([a, b] = c\) , \([a, c] = \beta b\) , \([b, c] = \gamma a\) , with \(\beta \neq 0\) and \(\gamma \neq 0\) , and g is simple. Let \(\phi\) be a canonical isomorphism of g onto the exterior power \(\bigwedge^{2}\mathfrak{g}^{*}\) of the dual space of \(\mathfrak{g}\) , determined to within an invertible constant factor (Algebra, Chapter III, § 8, no. 5); let \(u\) be the linear mapping of \(\mathfrak{g}\) into \(\mathfrak{g}^{*}\) defined by the condition \(\langle [x,y],u(z)\rangle = \langle x\wedge y,\phi (z)\rangle\) ; the bilinear form \(\Phi (x,y) = \langle x,u(y)\rangle\) on \(\mathfrak{g}\times \mathfrak{g}\) is symmetric and non-degenerate and \(2\Phi\) is the Killing form of \(\mathfrak{g}\) to within a factor \(\neq 0\) . For two 3-dimensional simple Lie algebras over K to be isomorphic, it is necessary and sufficient that the corresponding bilinear forms be equivalent to within a constant factor \(\neq 0\) .
\item
  If \(\mathbf{K}\) is not of characteristic 2, the simple algebra \(\mathfrak{g}\) defined in (c) is isomorphic to the Lie algebra \(\mathfrak{o}(\Phi)\) of the formal orthogonal group \(\mathbf{G}(\Phi)\) (§ 1, Exercise 26 (a)). For the existence in \(\mathfrak{g}\) of vectors \(x\) such that ad \(x\) admits eigenvectors not proportional to \(x\) , it is necessary and sufficient that \(\Phi\) be of index \(>0\) ; \(\mathfrak{g}\) then admits a basis \((a, b, c)\) such that \([a, b] = b\) , \([a, c] = -c\) , \([b, c] = a\) .
\item
  If K is of characteristic 2, show that there is no 2-mapping on g and therefore that g is not the Lie algebra of a formal group. Show that g admits derivations which are not inner.
\end{enumerate}

¶ 24. Let K be a field of arbitrary characteristic p.~We again adopt Definition 1.

\begin{enumerate}
\def\labelenumi{(\alph{enumi})}
\tightlist
\item
  Show that, unless \(p = n = 2\) , the only non-trivial ideals of \(\mathfrak{gl}(n, \mathbf{K})\)
\end{enumerate}

are the 1-dimensional centre \(\mathfrak{z}\) with basis \(\sum_{i=1}^{n} E_{ii}\) and the Lie algebra \(\mathfrak{sl}(n, \mathrm{K})\) consisting of the matrices of trace 0. (Except in the exceptional case indicated, note that if an ideal \(\mathfrak{a}\) contains one of the \(E_{ij}\) , it necessarily contains \(\mathfrak{sl}(n, \mathrm{K})\) . If \(\mathfrak{a}\) contains an element not belonging to \(\mathfrak{b}\) , by multiplying this element by at most four suitably chosen elements \(E_{ij}\) , a non-zero multiple of one of the \(E_{ij}\) is obtained.) If \(n\) is not a multiple of \(p\) , show that \(\mathfrak{gl}(n, \mathrm{K})\) is the direct sum of \(\mathfrak{z}\) and \(\mathfrak{sl}(n, \mathrm{K})\) and \(\mathfrak{sl}(n, \mathrm{K})\) is simple. If on the contrary \(n\) is a multiple of \(p\) and \(n > 2\) , show that \(\mathfrak{sl}(n, \mathrm{K}) / \mathfrak{z}\) is simple (same methods).

\begin{enumerate}
\def\labelenumi{(\alph{enumi})}
\setcounter{enumi}{1}
\tightlist
\item
  Show that, for \(n\) a multiple of \(p\) and \(n > 2\) , \(\mathfrak{gl}(n, \mathbf{K}) / \mathfrak{sl}(n, \mathbf{K})\) has radical \(\{0\}\) , but admits the quotient \(\mathfrak{gl}(n, \mathbf{K}) / \mathfrak{sl}(n, \mathbf{K})\) which is Abelian.
\end{enumerate}

¶ 25. Let K be a field of arbitrary characteristic p.~Let \(\mathfrak{sp}(2n, \mathrm{K})\) denote the Lie algebra of the formal symplectic group in 2n variables over K (§ 1, Exercise 26 (c)). Show that this Lie algebra, which is identified with a sub-algebra of \(\mathfrak{gl}(2n, \mathrm{K})\) , has basis consisting of the elements

\[
\begin{array}{c} H _ {i} = E _ {i i} - E _ {i + n, i + n} \quad (1 \leqslant i \leqslant n), \\ F _ {i j} = E _ {i j} - E _ {j + n, i + n} \quad (1 \leqslant i, j \leqslant n, i \neq j), \\ G _ {i j} ^ {\prime} = E _ {i, j + n} + E _ {j, i + n}, \qquad G _ {i j} ^ {\prime \prime} = E _ {i + n, j} + E _ {j + n, i} \quad (1 \leqslant i <   j \leqslant n), \\ E _ {i, i + n} \quad \text { and } \quad E _ {i + n, i} \quad (1 \leqslant i \leqslant n). \end{array}
\]

\begin{enumerate}
\def\labelenumi{(\alph{enumi})}
\item
  Show that, if \(p \neq 2\) and \(n \geqslant 1\) , the algebra \(\mathfrak{sp}(2n, \mathrm{K})\) is simple (for \(n = 1\) , \(\mathfrak{sp}(2, \mathrm{K}) = \mathfrak{sl}(2, \mathrm{K})\) ). (Same method as in Exercise 24.)
\item
  If \(p = 2\) and \(n \geqslant 3\) , show that the \(H_{i}, F_{ij}, G_{ij}'\) and \(G_{ij}''\) form an ideal \(\mathfrak{a}\) of dimension \(n(2n - 1)\) ; the elements of \(\mathfrak{a}\) of the form
\end{enumerate}

\[
\sum_ {i = 1} ^ {n} \lambda_ {i} H _ {i} + \sum_ {i \neq j} \alpha_ {i j} F _ {i j} + \sum_ {i <   j} \beta_ {i j} G _ {i j} ^ {\prime} + \sum_ {i <   j} \gamma_ {i j} G _ {i j} ^ {\prime \prime}
\]

such that \(\sum_{i=1}^{n} \lambda_i = 0\) form an ideal \(b \subset a\) of dimension \(n(2n - 1) - 1\) and the multiples of \(\sum_{i=1}^{n} H_i\) an ideal \(c\) , the centre of \(\mathfrak{sp}(2n, K)\) ; \(b, c\) are the only ideals of \(\mathfrak{a},\mathfrak{b} = \mathcal{D}(\mathfrak{sp}(2n,\mathbf{K}))\) and \(\mathfrak{b} / \mathfrak{c}\) is simple. (Same method.) What are the ideals of \(\mathfrak{sp}(4,\mathbf{K})\) when \(\mathbf{K}\) is of characteristic 2?

¶ 26. Let K be a field of characteristic ≠2 and Φ a non-degenerate symmetric bilinear form on an n-dimensional vector space over K.

\begin{enumerate}
\def\labelenumi{(\alph{enumi})}
\tightlist
\item
  Show that, if \(n \geqslant 5\) , the Lie algebra \(\mathfrak{o}(\Phi)\) is simple. (By extending the base field, it can be reduced to the case where the index of \(\Phi\) is \(m = [n/2]\) . If for example \(n = 2m\) is even, \(\mathfrak{o}(\Phi)\) has basis consisting of the elements
\end{enumerate}

\[
\begin{array}{c} H _ {i} = E _ {i i} - E _ {i + m, i + m}, \quad F _ {i j} = E _ {i j} - E _ {j + m, i + m} \quad (1 \leqslant i, j \leqslant m, i \neq j), \\ G _ {i j} ^ {\prime} = E _ {i, j + m} - E _ {j, i + m} \quad \text { and } \quad G _ {i j} ^ {\prime \prime} = E _ {i + m, j} - E _ {j + m, i} \quad (1 \leqslant i <   j \leqslant m); \end{array}
\]

then argue as in Exercises 24 and 25. Proceed similarly when \(n = 2m + 1\) is odd.)

\begin{enumerate}
\def\labelenumi{(\alph{enumi})}
\setcounter{enumi}{1}
\tightlist
\item
  Let \(\Delta\) be the discriminant of \(\Phi\) with respect to a basis. Suppose that \(n = 4\) . Show that, if \(\Delta\) is not a square in \(\mathbf{K}\) , \(\mathfrak{o}(\Phi)\) is simple. If \(\Delta\) is a square, \(\mathfrak{o}(\Phi)\) is the product of two isomorphic 3-dimensional simple Lie algebras. (Use the structure of \(\mathbf{G}(\Phi)\) described in Algebra, Chapter IX, §9, Exercise 16.) Deduce an example of a simple Lie algebra which becomes non-simple when extending the base field.
\end{enumerate}

\begin{enumerate}
\def\labelenumi{\arabic{enumi}.}
\setcounter{enumi}{26}
\tightlist
\item
  \begin{enumerate}
  \def\labelenumii{(\alph{enumii})}
  \tightlist
  \item
    Show that, in a finite-dimensional Lie \(p\) -algebra, the radical is a \(p\) -ideal. (Use Exercise 22 (c) of § 1.)
  \end{enumerate}
\end{enumerate}

\begin{enumerate}
\def\labelenumi{(\alph{enumi})}
\setcounter{enumi}{1}
\tightlist
\item
  For a Lie \(p\) -algebra \(\mathfrak{g}\) to have zero radical, it is necessary and sufficient that \(\mathfrak{g}\) contain no commutative \(p\) -ideal \(\neq \{0\}\) . (Use Exercise 22 (c) of § 1.)
\end{enumerate}

